\section{Benefit-Risk Assessment}
\label{sec:benefitrisk}
% ================================================================

The structured benefit-risk assessment below follows the FDA Benefit-Risk Framework
(PDUFA\,VI; \href{https://www.fda.gov/media/144657/download}{FDA Guidance, 2021}).

\vspace{3mm}

\begin{tcolorbox}[
  enhanced, breakable,
  colback=neutrallt, colframe=accent,
  title={\faBalanceScale\; Structured Benefit-Risk Summary --- VEN-ER 150\,mg},
  fonttitle=\bfseries\small\color{white},
  colbacktitle=accent,
  boxrule=0.7pt,
  left=8pt, right=8pt, top=6pt, bottom=6pt,
  attach boxed title to top left={yshift=-2mm, xshift=6mm}
]
\renewcommand{\arraystretch}{1.45}
\begin{tabularx}{\linewidth}{@{}>{\color{accent}\bfseries\small}l@{\hspace{8mm}}>{\raggedright\arraybackslash\small}X@{}}
\toprule
Dimension & Evidence \& Conclusion \\
\midrule
Analysis of Condition &
  MDD --- a severe, recurrent, disabling disorder with substantial unmet need for a
  flexible intermediate dosage strength to support titration. \\
Current Alternatives &
  Existing VEN-ER strengths (37.5, 75, 225\,mg); SSRIs; other SNRIs (duloxetine,
  desvenlafaxine). The 150\,mg addresses a clinically meaningful titration gap. \\
Benefits &
  Statistically significant and clinically meaningful MADRS CFB ($-$6.6 points vs.\ placebo,
  95\% CI $-$8.8 to $-$4.4; $p<0.0001$); response 62.1\% vs.\ 38.3\%; remission 41.6\%
  vs.\ 22.2\%; functional improvement (SDS $-$3.2; $p<0.0001$). \\
Risks &
  Nausea, dry mouth, dizziness --- class effects, predominantly mild--moderate, transient,
  manageable without dose reduction in most cases. Modest BP elevation
  ($+$3.1\,mmHg SBP). SAE rate comparable to reference product historical data. \\
Risk Management &
  Consistent with Effexor XR label (Black Box, Sections 5.1--5.11); no new safety signals;
  post-marketing cardiac monitoring commitment per RMP v2.1
  (\modref{m1/1-12-3/rmp.pdf}{m1/1.12.3/RMP-v2.1}). \\
Conclusion &
  \textbf{Favourable.} The benefit-risk profile of VEN-ER 150\,mg is positive. Clinical
  benefit in MDD is robust and consistent. Risks are known, manageable, and commensurate
  with the established SNRI class labelling. \\
\bottomrule
\end{tabularx}
\end{tcolorbox}

% ================================================================
